\documentclass[10pt,a4paper]{amsart}
\usepackage{amsmath,amssymb,amsthm,graphicx}
% Retained minimal preamble for the Stacks Project chapter more-morphisms.tex (More on Morphisms).
% Verbatim excerpt of the chapter preamble: Xy-pic setup (native preamble lines 14-18) and the native
% theorem declarations with the mathematical macros (native preamble lines 154-217). Copied mechanically;
% nothing added, removed or redefined. The wrapper class and the amsmath/amssymb/amsthm/graphicx base
% packages are supplied by the assembler, so the class/externaldocument part of the native preamble is omitted.
% For commutative diagrams we use Xy-pic
\usepackage[all]{xy}

% We use 2cell for 2-commutative diagrams.
\xyoption{2cell}
% Theorem environments.
%
\theoremstyle{plain}
\newtheorem{theorem}[subsection]{Theorem}
\newtheorem{proposition}[subsection]{Proposition}
\newtheorem{lemma}[subsection]{Lemma}

\theoremstyle{definition}
\newtheorem{definition}[subsection]{Definition}
\newtheorem{example}[subsection]{Example}
\newtheorem{exercise}[subsection]{Exercise}
\newtheorem{situation}[subsection]{Situation}

\theoremstyle{remark}
\newtheorem{remark}[subsection]{Remark}
\newtheorem{remarks}[subsection]{Remarks}

\numberwithin{equation}{subsection}

% Macros
%
\def\lim{\mathop{\mathrm{lim}}\nolimits}
\def\colim{\mathop{\mathrm{colim}}\nolimits}
\def\Spec{\mathop{\mathrm{Spec}}}
\def\Hom{\mathop{\mathrm{Hom}}\nolimits}
\def\Ext{\mathop{\mathrm{Ext}}\nolimits}
\def\SheafHom{\mathop{\mathcal{H}\!\mathit{om}}\nolimits}
\def\SheafExt{\mathop{\mathcal{E}\!\mathit{xt}}\nolimits}
\def\Sch{\mathit{Sch}}
\def\Mor{\mathop{\mathrm{Mor}}\nolimits}
\def\Ob{\mathop{\mathrm{Ob}}\nolimits}
\def\Sh{\mathop{\mathit{Sh}}\nolimits}
\def\NL{\mathop{N\!L}\nolimits}
\def\CH{\mathop{\mathrm{CH}}\nolimits}
\def\proetale{{pro\text{-}\acute{e}tale}}
\def\etale{{\acute{e}tale}}
\def\QCoh{\mathit{QCoh}}
\def\Ker{\mathop{\mathrm{Ker}}}
\def\Im{\mathop{\mathrm{Im}}}
\def\Coker{\mathop{\mathrm{Coker}}}
\def\Coim{\mathop{\mathrm{Coim}}}

% Boxtimes
%
\DeclareMathSymbol{\boxtimes}{\mathbin}{AMSa}{"02}

%
% Macros for moduli stacks/spaces
%
\def\QCohstack{\mathcal{QC}\!\mathit{oh}}
\def\Cohstack{\mathcal{C}\!\mathit{oh}}
\def\Spacesstack{\mathcal{S}\!\mathit{paces}}
\def\Quotfunctor{\mathrm{Quot}}
\def\Hilbfunctor{\mathrm{Hilb}}
\def\Curvesstack{\mathcal{C}\!\mathit{urves}}
\def\Polarizedstack{\mathcal{P}\!\mathit{olarized}}
\def\Complexesstack{\mathcal{C}\!\mathit{omplexes}}
% \Pic is the operator that assigns to X its picard group, usage \Pic(X)
% \Picardstack_{X/B} denotes the Picard stack of X over B
% \Picardfunctor_{X/B} denotes the Picard functor of X over B
\def\Pic{\mathop{\mathrm{Pic}}\nolimits}
\def\Picardstack{\mathcal{P}\!\mathit{ic}}
\def\Picardfunctor{\mathrm{Pic}}
\def\Deformationcategory{\mathcal{D}\!\mathit{ef}}
\usepackage{cleveref}
\crefname{originalreference}{original reference}{original references}
\Crefname{originalreference}{Original reference}{Original references}
\makeatletter
\newlabel{algebra-lemma-dimension-at-a-point-preserved-field-extension}{{00P4}{source}{}{sourcepointer-1}{}}
\newlabel{algebra-lemma-dimension-at-a-point-preserved-field-extension@cref}{{[originalreference][0][]00P4}{[1][0][]source}{}{}{}}
\makeatother
\crefname{originalreference}{original reference}{original references}
\Crefname{originalreference}{Original reference}{Original references}
\makeatletter
\newlabel{algebra-lemma-dominate}{{00IA}{source}{}{sourcepointer-2}{}}
\newlabel{algebra-lemma-dominate@cref}{{[originalreference][0][]00IA}{[1][0][]source}{}{}{}}
\makeatother
\crefname{originalreference}{original reference}{original references}
\Crefname{originalreference}{Original reference}{Original references}
\makeatletter
\newlabel{algebra-lemma-finite-type-domain-over-valuation-ring-dim-fibres}{{00QK}{source}{}{sourcepointer-3}{}}
\newlabel{algebra-lemma-finite-type-domain-over-valuation-ring-dim-fibres@cref}{{[originalreference][0][]00QK}{[1][0][]source}{}{}{}}
\makeatother
\crefname{originalreference}{original reference}{original references}
\Crefname{originalreference}{Original reference}{Original references}
\makeatletter
\newlabel{morphisms-lemma-dimension-fibre-at-a-point}{{02FX}{source}{}{sourcepointer-4}{}}
\newlabel{morphisms-lemma-dimension-fibre-at-a-point@cref}{{[originalreference][0][]02FX}{[1][0][]source}{}{}{}}
\makeatother
\title{Fibre dimension plus transcendence degree equals the transcendence degree of a flat local extension}
\author{The Stacks Project Authors}
\date{stacks-project:more-morphisms.tex:local-snapshot-20261011; GNU FDL 1.2 or later}
\begin{document}
\maketitle
\noindent\textit{Editorial note.} Counted span is the author's own proof environment, source lines 14770-14814, immediately following the statement at 14754-14768.\par\medskip
\noindent\textit{Editorial note.} Four cross-chapter pointers (Algebra, Morphisms) are recorded in assembly.external\_references.\par\medskip
\noindent\textit{Editorial note.} The two reduction steps (finite type model, valuation ring case) are both displayed inside the span.\par\medskip
\noindent\textit{Allowed standard prerequisites (editorial).}\begin{itemize}
\item Dimension of the fibre at a point for a flat local extension. Morphisms, lemma morphisms-lemma-dimension-fibre-at-a-point, cited in the proof. Identifies the fibre dimension with the dimension of the special fibre of the local extension.
\item Dimension preserved under field extension. Algebra, lemma algebra-lemma-dimension-at-a-point-preserved-field-extension, cited in the proof. Keeps the fibre dimension unchanged when passing to a suitable localisation.
\item Finite type domains over a valuation ring and their fibres. Algebra, lemma algebra-lemma-finite-type-domain-over-valuation-ring-dim-fibres, cited at proof lines 14811-14812. Gives the equality of dimension of special fibre and generic fibre in the valuation ring case that the proof reduces to.
\item Dominating a local map by a finite type map of domains. Algebra, lemma algebra-lemma-dominate, cited in the proof. Replaces B by a domain C of finite type over A (with B a localisation of C), which is what allows the reduction to valuation rings.
\end{itemize}\medskip
\ifdefined\hypertarget\hypertarget{sourcepointer-1}{}\fi
\noindent\textit{Source reference.} Original reference 00P4: focused-seeds/geometry/S9-morph-equality-dimensions/ref-tags.txt, source lines 4--4. Navigation pointer only: the counted proof cites this statement by \textbackslash{}ref as a standing input; the deduction the counted proof has to make is carried out inside the counted span, and no part of that deduction is delegated to this pointer.\par\medskip
\ifdefined\hypertarget\hypertarget{sourcepointer-2}{}\fi
\noindent\textit{Source reference.} Original reference 00IA: focused-seeds/geometry/S9-morph-equality-dimensions/ref-tags.txt, source lines 5--5. Navigation pointer only: the counted proof cites this statement by \textbackslash{}ref as a standing input; the deduction the counted proof has to make is carried out inside the counted span, and no part of that deduction is delegated to this pointer.\par\medskip
\ifdefined\hypertarget\hypertarget{sourcepointer-3}{}\fi
\noindent\textit{Source reference.} Original reference 00QK: focused-seeds/geometry/S9-morph-equality-dimensions/ref-tags.txt, source lines 6--6. Navigation pointer only: the counted proof cites this statement by \textbackslash{}ref as a standing input; the deduction the counted proof has to make is carried out inside the counted span, and no part of that deduction is delegated to this pointer.\par\medskip
\ifdefined\hypertarget\hypertarget{sourcepointer-4}{}\fi
\noindent\textit{Source reference.} Original reference 02FX: focused-seeds/geometry/S9-morph-equality-dimensions/ref-tags.txt, source lines 7--7. Navigation pointer only: the counted proof cites this statement by \textbackslash{}ref as a standing input; the deduction the counted proof has to make is carried out inside the counted span, and no part of that deduction is delegated to this pointer.\par\medskip
\begin{lemma}
\label{lemma-equality-dimensions}
Let $A \to B$ be a local homomorphism of local rings. Assume
\begin{enumerate}
\item $A$ and $B$ are domains and $A \subset B$,
\item $B$ is essentially of finite type over $A$, and
\item $B$ is flat over $A$.
\end{enumerate}
Then we have
$$
\dim(B/\mathfrak m_AB) +
\text{trdeg}_{\kappa(\mathfrak m_A)}(\kappa(\mathfrak m_B)) =
\text{trdeg}_A(B).
$$
\end{lemma}

\begin{proof}
Let $C$ be an $A$-algebra of finite type and $\mathfrak q$ a prime of $C$
such that $B = C_{\mathfrak q}$. We may assume $C$ is a domain.
We have
$\dim_{\mathfrak q}(C/\mathfrak m_AC) = \dim(B/\mathfrak m_AB) +
\text{trdeg}_{\kappa(\mathfrak m_A)}(\kappa(\mathfrak m_B))$, see
Morphisms, Lemma \ref{morphisms-lemma-dimension-fibre-at-a-point}.
Setting $K$ equal to the fraction field of $A$
we see by the same reference that
$\dim(C \otimes_A K) = \text{trdeg}_A(B)$.
Thus we are really trying to prove that
$\dim_{\mathfrak q}(C/\mathfrak m_AC) = \dim(C \otimes_A K)$.
Choose a valuation ring $A'$ in $K$ dominating $A$, see
Algebra, Lemma \ref{algebra-lemma-dominate}.
Set $C' = C \otimes_A A'$.
Choose a prime $\mathfrak q'$ of $C'$ lying over $\mathfrak q$; such a
prime exists because
$$
C'/\mathfrak m_{A'}C' =
C/\mathfrak m_AC \otimes_{\kappa(\mathfrak m_A)} \kappa(\mathfrak m_{A'})
$$
which proves that $C/\mathfrak m_AC \to C'/\mathfrak m_{A'}C'$ is faithfully
flat. This also proves that
$\dim_{\mathfrak q}(C/\mathfrak m_AC) =
\dim_{\mathfrak q'}(C'/\mathfrak m_{A'}C')$, see
Algebra,
Lemma \ref{algebra-lemma-dimension-at-a-point-preserved-field-extension}.
Note that $B' = C'_{\mathfrak q'}$ is a localization of $B \otimes_A A'$.
Hence $B'$ is flat over $A'$. The generic fibre $B' \otimes_{A'} K$
is a localization of $B \otimes_A K$. Hence $B'$ is a domain.
If we prove the lemma for $A' \subset B'$, then we get the equality
$\dim_{\mathfrak q'}(C'/\mathfrak m_{A'}C') = \dim(C' \otimes_{A'} K)$
which implies the desired equality
$\dim_{\mathfrak q}(C/\mathfrak m_AC) = \dim(C \otimes_A K)$
by what was said above. This reduces the
lemma to the case where $A$ is a valuation ring.

\medskip\noindent
Let $A \subset B$ be as in the lemma with $A$ a valuation ring.
As before write $B = C_{\mathfrak q}$ for some domain $C$ of finite
type over $A$. By
Algebra,
Lemma \ref{algebra-lemma-finite-type-domain-over-valuation-ring-dim-fibres}
we obtain $\dim(C/\mathfrak m_AC) = \dim(C \otimes_A K)$ and we win.
\end{proof}

\end{document}
